\documentclass{amsart}
% For dealing with references we use the comment environment
\usepackage{verbatim}
\newenvironment{reference}{\comment}{\endcomment}
%\newenvironment{reference}{}{}
\newenvironment{slogan}{\comment}{\endcomment}
\newenvironment{history}{\comment}{\endcomment}

% For commutative diagrams we use Xy-pic
\usepackage[all]{xy}

% We use 2cell for 2-commutative diagrams.
\xyoption{2cell}
\UseAllTwocells

% We use multicol for the list of chapters between chapters
\usepackage{multicol}

% This is generally recommended for better output
\usepackage{lmodern}
\usepackage[T1]{fontenc}

% For cross-file-references

% Package for hypertext links:
\usepackage{hyperref}

% For any local file, say "hello.tex" you want to link to please

% Theorem environments.
%
\theoremstyle{plain}
\newtheorem{theorem}[subsection]{Theorem}
\newtheorem{proposition}[subsection]{Proposition}
\newtheorem{lemma}[subsection]{Lemma}

\theoremstyle{definition}
\newtheorem{definition}[subsection]{Definition}
\newtheorem{example}[subsection]{Example}
\newtheorem{exercise}[subsection]{Exercise}
\newtheorem{situation}[subsection]{Situation}

\theoremstyle{remark}
\newtheorem{remark}[subsection]{Remark}
\newtheorem{remarks}[subsection]{Remarks}

\numberwithin{equation}{subsection}

% Macros
%
\def\lim{\mathop{\mathrm{lim}}\nolimits}
\def\colim{\mathop{\mathrm{colim}}\nolimits}
\def\Spec{\mathop{\mathrm{Spec}}}
\def\Hom{\mathop{\mathrm{Hom}}\nolimits}
\def\Ext{\mathop{\mathrm{Ext}}\nolimits}
\def\SheafHom{\mathop{\mathcal{H}\!\mathit{om}}\nolimits}
\def\SheafExt{\mathop{\mathcal{E}\!\mathit{xt}}\nolimits}
\def\Sch{\mathit{Sch}}
\def\Mor{\mathop{\mathrm{Mor}}\nolimits}
\def\Ob{\mathop{\mathrm{Ob}}\nolimits}
\def\Sh{\mathop{\mathit{Sh}}\nolimits}
\def\NL{\mathop{N\!L}\nolimits}
\def\CH{\mathop{\mathrm{CH}}\nolimits}
\def\proetale{{pro\text{-}\acute{e}tale}}
\def\etale{{\acute{e}tale}}
\def\QCoh{\mathit{QCoh}}
\def\Ker{\mathop{\mathrm{Ker}}}
\def\Im{\mathop{\mathrm{Im}}}
\def\Coker{\mathop{\mathrm{Coker}}}
\def\Coim{\mathop{\mathrm{Coim}}}

% Boxtimes
%
\DeclareMathSymbol{\boxtimes}{\mathbin}{AMSa}{"02}

%
% Macros for moduli stacks/spaces
%
\def\QCohstack{\mathcal{QC}\!\mathit{oh}}
\def\Cohstack{\mathcal{C}\!\mathit{oh}}
\def\Spacesstack{\mathcal{S}\!\mathit{paces}}
\def\Quotfunctor{\mathrm{Quot}}
\def\Hilbfunctor{\mathrm{Hilb}}
\def\Curvesstack{\mathcal{C}\!\mathit{urves}}
\def\Polarizedstack{\mathcal{P}\!\mathit{olarized}}
\def\Complexesstack{\mathcal{C}\!\mathit{omplexes}}
% \Pic is the operator that assigns to X its picard group, usage \Pic(X)
% \Picardstack_{X/B} denotes the Picard stack of X over B
% \Picardfunctor_{X/B} denotes the Picard functor of X over B
\def\Pic{\mathop{\mathrm{Pic}}\nolimits}
\def\Picardstack{\mathcal{P}\!\mathit{ic}}
\def\Picardfunctor{\mathrm{Pic}}
\def\Deformationcategory{\mathcal{D}\!\mathit{ef}}

\setlength{\emergencystretch}{3em}
\begin{document}
\title[Stacks Project, Tag 03SC]{284 --- Vanishing, finiteness, and algebraically closed base change for torsion etale cohomology of curves}
\maketitle
\noindent Copyright (C) 2005--2025 Johan de Jong. Stacks Project authors. Source: \href{https://stacks.math.columbia.edu/tag/03SC}{Tag 03SC}.

\noindent Editorial difficulty note (not author text). Difficulty H2: Arbitrary torsion coefficients must be reduced to constructible prime-field local systems without losing any of eight conclusions. The proof performs prime-power devissage and separates a dense open into irreducible components, then uses extension by zero to assemble the curve statements from the established local-system case..

\noindent Editorial provenance note (not author text). History: excerpt from revision \texttt{a0444\allowbreak 6e57e\allowbreak c1fbc\allowbreak 252a8\allowbreak 71afc\allowbreak ec775\allowbreak 2fb28\allowbreak 07b14}, etale-cohomology.tex, lines 15655--15737. Reference presentation was adapted; original-author context and core support blocks were retained or added. The mathematical wording is unchanged. Licensed under GNU FDL 1.2 or later; no invariant sections or cover texts. See ../licenses/COPYING. Context blocks reproduce author definitions and supporting results; they are not additional counted problems.

\begin{situation}
\label{situation-what-to-prove}
Here $k$ is an algebraically closed field, $X$ is a separated, finite type
scheme of dimension $\leq 1$ over $k$, and $\mathcal{F}$ is a torsion
abelian sheaf on $X_\etale$.
\end{situation}

\noindent
In Situation \ref{situation-what-to-prove}
we want to prove the following statements
\begin{enumerate}
\item
\label{item-vanishing}
$H^q_\etale(X, \mathcal{F}) = 0$ for $q > 2$,
\item
\label{item-vanishing-affine}
$H^q_\etale(X, \mathcal{F}) = 0$ for $q > 1$ if $X$ is affine,
\item
\label{item-vanishing-p-p}
$H^q_\etale(X, \mathcal{F}) = 0$ for $q > 1$ if $p = \text{char}(k) > 0$
and $\mathcal{F}$ is $p$-power torsion,
\item
\label{item-finite-prime-to-p}
$H^q_\etale(X, \mathcal{F})$ is finite if $\mathcal{F}$ is
constructible and torsion prime to $\text{char}(k)$,
\item
\label{item-finite-proper}
$H^q_\etale(X, \mathcal{F})$ is finite if $X$ is proper and
$\mathcal{F}$ constructible,
\item
\label{item-base-change-prime-to-p}
$H^q_\etale(X, \mathcal{F}) \to
H^q_\etale(X_{k'}, \mathcal{F}|_{X_{k'}})$ is an isomorphism
for any extension $k'/k$ of algebraically closed fields
if $\mathcal{F}$ is torsion prime to $\text{char}(k)$,
\item
\label{item-base-change-proper}
$H^q_\etale(X, \mathcal{F}) \to
H^q_\etale(X_{k'}, \mathcal{F}|_{X_{k'}})$ is an isomorphism
for any extension $k'/k$ of algebraically closed fields
if $X$ is proper,
\item
\label{item-surjective}
$H^2_\etale(X, \mathcal{F}) \to H^2_\etale(U, \mathcal{F})$
is surjective for all $U \subset X$ open.
\end{enumerate}
Given any Situation \ref{situation-what-to-prove}
we will say that
``statements (\ref{item-vanishing}) -- (\ref{item-surjective}) hold''
if those statements that apply to the given situation are true.


\begin{theorem}
\label{theorem-vanishing-affine-curves}
If $k$ is an algebraically closed field, $X$ is a separated, finite type
scheme of dimension $\leq 1$ over $k$, and $\mathcal{F}$ is a torsion
abelian sheaf on $X_\etale$, then
\begin{enumerate}
\item
$H^q_\etale(X, \mathcal{F}) = 0$ for $q > 2$,
\item
$H^q_\etale(X, \mathcal{F}) = 0$ for $q > 1$ if $X$ is affine,
\item
$H^q_\etale(X, \mathcal{F}) = 0$ for $q > 1$ if $p = \text{char}(k) > 0$
and $\mathcal{F}$ is $p$-power torsion,
\item
$H^q_\etale(X, \mathcal{F})$ is finite if $\mathcal{F}$ is
constructible and torsion prime to $\text{char}(k)$,
\item
$H^q_\etale(X, \mathcal{F})$ is finite if $X$ is proper and
$\mathcal{F}$ constructible,
\item
$H^q_\etale(X, \mathcal{F}) \to
H^q_\etale(X_{k'}, \mathcal{F}|_{X_{k'}})$ is an isomorphism
for any extension $k'/k$ of algebraically closed fields
if $\mathcal{F}$ is torsion prime to $\text{char}(k)$,
\item
$H^q_\etale(X, \mathcal{F}) \to
H^q_\etale(X_{k'}, \mathcal{F}|_{X_{k'}})$ is an isomorphism
for any extension $k'/k$ of algebraically closed fields
if $X$ is proper,
\item
$H^2_\etale(X, \mathcal{F}) \to H^2_\etale(U, \mathcal{F})$
is surjective for all $U \subset X$ open.
\end{enumerate}
\end{theorem}

\begin{proof}
The theorem says that in Situation \ref{situation-what-to-prove}
statements (\ref{item-vanishing}) -- (\ref{item-surjective}) hold.
Our first step is to replace $X$ by its reduction, which is permissible
by Proposition \href{https://stacks.math.columbia.edu/tag/03SI}{proposition topological invariance (Tag 03SI)}.
By Lemma \href{https://stacks.math.columbia.edu/tag/03SA}{lemma torsion colimit constructible (Tag 03SA)} we can write
$\mathcal{F}$ as a filtered colimit of constructible abelian sheaves.
Taking cohomology commutes with colimits, see Lemma \href{https://stacks.math.columbia.edu/tag/03Q5}{lemma colimit (Tag 03Q5)}.
Moreover, pullback via $X_{k'} \to X$ commutes with colimits as a left
adjoint. Thus it suffices to prove the statements for a constructible sheaf.

\medskip\noindent
In this paragraph we use Lemma \href{https://stacks.math.columbia.edu/tag/0A5C}{lemma ses statements (Tag 0A5C)} without further
mention. Writing
$\mathcal{F} = \mathcal{F}_1 \oplus \ldots \oplus \mathcal{F}_r$
where $\mathcal{F}_i$ is $\ell_i$-primary for some prime $\ell_i$, we may
assume that $\ell^n$ kills $\mathcal{F}$ for some prime $\ell$. Now consider
the exact sequence
$$
0 \to \mathcal{F}[\ell] \to \mathcal{F} \to \mathcal{F}/\mathcal{F}[\ell] \to 0.
$$
Thus we see that it suffices to assume that $\mathcal{F}$ is $\ell$-torsion.
This means that $\mathcal{F}$ is a constructible sheaf of
$\mathbf{F}_\ell$-vector spaces for some prime number $\ell$.

\medskip\noindent
By definition this means there is a dense open $U \subset X$
such that $\mathcal{F}|_U$ is finite locally constant sheaf of
$\mathbf{F}_\ell$-vector spaces. Since $\dim(X) \leq 1$ we may
assume, after shrinking $U$, that $U = U_1 \amalg \ldots \amalg U_n$
is a disjoint union of irreducible schemes (just remove the closed
points which lie in the intersections of $\geq 2$ components of $U$).
By Lemma \href{https://stacks.math.columbia.edu/tag/0GJA}{lemma restrict to open (Tag 0GJA)} we reduce to the case
$\mathcal{F} = j_!\mathcal{G}$ where $\mathcal{G}$ is a finite
locally constant sheaf of $\mathbf{F}_\ell$-vector spaces on $U$.

\medskip\noindent
Since we chose $U = U_1 \amalg \ldots \amalg U_n$ with $U_i$ irreducible
we have
$$
j_!\mathcal{G} =
j_{1!}(\mathcal{G}|_{U_1}) \oplus \ldots \oplus
j_{n!}(\mathcal{G}|_{U_n})
$$
where $j_i : U_i \to X$ is the inclusion morphism.
The case of $j_{i!}(\mathcal{G}|_{U_i})$ is handled in
Lemma \href{https://stacks.math.columbia.edu/tag/03SD}{lemma vanishing easier (Tag 03SD)}.
\end{proof}
\end{document}
